\section{Randomized bi-Lipschitz pushforward and W$_2$ covariance
envelope}\label{sec:methodology}

\subsection{Theoretical framework}

The framework converts distance between distribution-valued firm
characteristics into a conditional restriction on systematic covariance.
To make that argument transparent, the construction separates
characteristic-side information, the mechanism that transmits characteristics
into exposures, and the joint arrangement of the resulting exposure draws.
For firm $i$, the law $C_i$ records a distribution over risk-relevant
locations rather than collapsing heterogeneous information to one vector.
A common map sends draws from these laws to factor exposures, while a
Lipschitz restriction controls how much the map can distort characteristic distance.
The map, its distortion and slack bounds, and exposure magnitudes are
maintained rather than estimated.

Because the restriction concerns systematic covariance, the return
representation is written to accommodate either a finite or an infinite
collection of common risk directions.
A real separable Hilbert space generalizes the familiar Euclidean loading
space: its inner product measures alignment and its norm measures magnitude.
Let $\mathcal H$ be such a space, let $F_t\in\mathcal H$ be a centered factor
innovation with covariance operator $\Gamma$, let $e_t^i$ be a centered
idiosyncratic component, and write $\widetilde r_t^i$ for firm $i$'s centered return:
\begin{equation}\label{eq:hilbert-return}
  \widetilde r_t^i = \inner{B_i,F_t}_{\mathcal H}+e_t^i.
\end{equation}
The return equation contains two familiar sources of randomness: $F_t$
represents shocks common to firms, whereas $e_t^i$ represents the centered
firm-specific return component.
The random element $B_i$ adds a distinct source of variation by allowing a
firm's loading on common shocks to follow an exposure law.
To connect that loading law to distribution-valued characteristics, let
$B_i=T(X_i,U_i)$, where $X_i\sim C_i$ is a characteristic draw and $U_i$
allows the same characteristic to map to different exposures.
Thus $X_i$ represents heterogeneity within the firm's characteristic law,
while $U_i$ represents transmission variation conditional on a
characteristic; both differ from the return innovations $F_t$ and $e_t^i$.
Equivalently, a stochastic kernel $K(x,\cdot)=\mathcal L(T(x,U))$ maps
characteristics to exposure laws and $P_i=C_i\,K=\mathcal L(B_i)$.
The deterministic map and additive-slack cases are special cases of this construction.

Transport comparisons require the firm laws to share a characteristic domain
whose metric supplies the cost of moving between risk-relevant locations.
Let $(\Omega,d_\Omega)$ be a metric space of semantic risk locations with
Borel $\sigma$-algebra $\mathcal B$.
Fix any $\omega_0\in\Omega$ and define the admissible characteristic domain
\begin{equation}\label{eq:characteristic-domain}
  \mathcal P_1(\Omega)
  :=\left\{C:\ C\text{ is a Borel probability measure on }\Omega,\quad
  \int_\Omega d_\Omega(\omega,\omega_0)\,C(d\omega)<\infty\right\}.
\end{equation}
The definition is independent of the chosen base point by the triangle inequality.

The object to be bounded is the covariance of systematic components, so
exposures are first expressed in the factor-risk geometry induced by $\Gamma$.
Under a coupling $\pi$ of two exposure laws, define
\begin{equation}\label{eq:hilbert-systematic-covariance}
  \kappa_\pi=\E_\pi\inner{Z_i,Z_j}_{\mathcal H},
  \qquad Z_i:=\Gamma^{1/2}B_i.
\end{equation}
Thus $\Gamma$ weights factor directions by their contribution to systematic
risk, and $Z_i$ measures exposure in the covariance-relevant geometry.
Unlike $X_i$ or $U_i$, the coupling $\pi$ is not a further marginal shock; it
specifies how draws from two already-defined exposure laws occur jointly.

The envelope itself concerns systematic covariance, whereas its
interpretation using total returns requires a separate return bridge.
Under that bridge, the residual is square-integrable and centered, is
orthogonal to $F_t$, and has zero cross-firm residual covariance under the
same coupling $\pi$.
The full exposure laws $P_i$ are the objects entering the transport envelope.
The article cloud is only a finite-support proxy for $C_i$, and the data do
not identify the latent draws, the kernel, or the joint coupling of the exposure laws.

\subsection{Randomized bi-Lipschitz pushforward and the W$_2$ covariance
envelope}\label{sec:random-exposure}

Characteristic distance can restrict covariance only if differences in
characteristics survive their transmission into systematic exposures.
The formulation therefore treats the characteristic-to-exposure link as a
randomized pushforward: a common carrier captures the shared part of the
link, while the kernel draw allows exposure heterogeneity around it.
This structure also isolates the object later linked to returns: the excess
of exposure distance under a realized coupling over the least possible
transport cost of the latent exposure laws.

Quadratic transport and covariance both require second moments, so this
subsection strengthens the characteristic domain accordingly.
Let $\Omega$ be a Polish metric space with its Borel $\sigma$-algebra, let
$\mathcal H$ be the exposure Hilbert space, and let $\Gamma:\mathcal
H\to\mathcal H$ be a bounded, positive self-adjoint operator.
Its unique bounded positive square root is denoted by $\Gamma^{1/2}$.
Write
\begin{equation}\label{eq:random-p2-domain}
  \mathcal P_2(\Omega)
  :=\left\{C\in\mathcal P_1(\Omega):
  \int_\Omega d_\Omega(\omega,\omega_0)^2\,C(d\omega)<\infty\right\}.
\end{equation}
The lower transfer bound will need to pull a coupling on the carrier image
back to characteristic space, which motivates the inverse condition stated next.
The common carrier $t:\Omega\to\mathcal H$ is Borel measurable, and
$t_\Gamma:=\Gamma^{1/2}\circ t$ is a Borel embedding with a Borel inverse on its image.
The stochastic map $T:\Omega\times\mathcal U\to\mathcal H$ is Borel
measurable, and the maintained integrability conditions put the induced
exposure laws in $\mathcal P_2(\mathcal H)$.

The assumption below allows two controlled departures from exact common transmission.
Economically, $L$ measures how faithfully characteristic differences survive
as systematic-exposure differences: $L=1$ preserves distances, whereas $L>1$
permits multiplicative compression or expansion.
The bound $\tau_i$ permits an additive firm-level departure from the common
carrier in the same risk units.
This slack is a transmission tolerance rather than article-measurement error,
and the kernel draw $U_i$ that generates it remains distinct from the
cross-firm coupling $\pi$ used to join the firm-level exposure laws.
\begin{assumption}[Common bi-Lipschitz random map]
  \label{ass:random-common-map}
  Let $C_i,C_j\in\mathcal P_2(\Omega)$ and let $t:\Omega\to\mathcal H$ be
  the common carrier. Assume that, for some $L\ge1$,
  \begin{equation}\label{eq:random-bilipschitz}
    L^{-1}d_\Omega(x,y)\le\norm{t_\Gamma(x)-t_\Gamma(y)}_{\mathcal H}
    \le Ld_\Omega(x,y).
  \end{equation}
  If $X_i\sim C_i$ and $U_i$ is drawn from the common transmission kernel,
  define the exposure random element
  \begin{equation}\label{eq:random-exposure-pushforward}
    B_i=T(X_i,U_i)
    =t(X_i)+\varepsilon_i(X_i,U_i),
    \qquad \norm{\Gamma^{1/2}\varepsilon_i(x,u)}_{\mathcal H}\le\tau_i,
    \qquad P_i:=\mathcal L(B_i).
  \end{equation}
  The slack maps $\varepsilon_i$ are Borel measurable, and the preceding
  integrability condition ensures that $P_i$ has a finite
  $\Gamma$-weighted second moment. The map, the common Lipschitz scale $L$,
  and the deviation bounds $\tau_i$ are maintained over the analyzed regime.
  Both scales are stated in risk coordinates:
  equation~\ref{eq:random-bilipschitz}
  constrains $\Gamma^{1/2}t$ rather than $t$ alone, and $\tau_i$ bounds the
  deviation after the same factor weighting. Every transport quantity
  displayed below therefore lives in the one geometry that the covariance
  envelope is defined in, and no step silently exchanges
  $\norm{\cdot}_{\mathcal H}$ for $\norm{\Gamma^{1/2}\cdot}_{\mathcal H}$.
\end{assumption}

The common-map restriction does the economic work of connecting
characteristic and exposure geometry; once the exposure laws are fixed, the
covariance envelope follows from an elementary inner-product identity.
We state that identity first because every bound below is an instance of it.

\begin{lemma}[Polarization Identity]\label{lem:polarization}
  For any $x,y$ in a real inner-product space,
  \begin{equation}\label{eq:polarization}
    \inner{x,y}
    =\frac12\left(\norm{x}^2+\norm{y}^2-\norm{x-y}^2\right).
  \end{equation}
\end{lemma}

Polarization supplies the specific link between quadratic transport and
covariance: its squared ground cost is exactly the quantity that
equation~\ref{eq:polarization} converts into an inner product.
Minimizing squared transport cost over couplings therefore maximizes
covariance over couplings, which gives W$_2$ a theoretical role not shared by
the alternative distances used later as empirical comparisons.

To determine what the marginal exposure laws alone can say about covariance,
we must compare their possible joint arrangements.
A coupling is a joint probability law with the prescribed firm-level
marginals; economically, it records which exposure draws occur together while
leaving each firm's marginal exposure law unchanged.
W$_2$ searches across these dependence arrangements for the smallest expected
squared exposure distance.

Three quantities summarize the pair in factor-risk coordinates: the marginal
second moment $v_i$, the coupling-specific covariance $\kappa_\pi$, and the
minimum squared exposure distance $W_{2,\Gamma}^2$.
For $Z_i=\Gamma^{1/2}B_i$, write
\begin{equation}\label{eq:random-covariance-objects}
  v_i:=\E\norm{Z_i}^{2},\qquad
  \kappa_\pi:=\E_\pi\inner{Z_i,Z_j},\qquad
  W_{2,\Gamma}^{2}(P_i,P_j):=\inf_\pi
  \E_\pi\norm{Z_i-Z_j}^{2}.
\end{equation}
The infimum ranges over couplings of $P_i$ and $P_j$ after the covariance
geometry induced by $\Gamma$.
Because the marginal second moments remain fixed as the coupling changes,
define the candidate upper endpoint
\begin{equation}\label{eq:random-covariance-envelope}
  \bar\kappa_{ij}:=\frac12\left(v_i+v_j-W_{2,\Gamma}^{2}(P_i,P_j)\right).
\end{equation}

The central result establishes that this endpoint is the attainable
covariance ceiling and identifies the gap between that ceiling and any
particular coupling.

\begin{theorem}[Random-exposure covariance envelope]
  \label{thm:random-exposure-envelope}
  Under \Cref{ass:random-common-map}, for every coupling $\pi$ of the exposure
  laws,
  \begin{equation}\label{eq:random-exposure-gap}
    \kappa_\pi
    =\bar\kappa_{ij}-\frac12\Delta_\pi,
    \qquad
    \Delta_\pi:=\E_\pi\norm{Z_i-Z_j}^{2}
    -W_{2,\Gamma}^{2}(P_i,P_j).
  \end{equation}
  The excess is nonnegative,
  \begin{equation}\label{eq:random-exposure-excess-nonneg}
    \Delta_\pi\ge0 ,
  \end{equation}
  and consequently $\bar\kappa_{ij}$ is the greatest covariance attainable by
  a coupling of the two exposure laws, with every realized coupling satisfying
  $\kappa_\pi\le\bar\kappa_{ij}$.
\end{theorem}

\begin{proof}
  Polarization in $\mathcal H$ gives
  $\norm{Z_i-Z_j}^2=\norm{Z_i}^2+\norm{Z_j}^2-2\inner{Z_i,Z_j}$ pointwise.
  Integrating against $\pi$ and using that its marginals are $P_i$ and $P_j$
  replaces the first two terms by $v_i$ and $v_j$, which depend on the
  marginals alone, so $\E_\pi\norm{Z_i-Z_j}^2=v_i+v_j-2\kappa_\pi$.
  Subtracting $W_{2,\Gamma}^2(P_i,P_j)$ and dividing by two gives
  equation~\ref{eq:random-exposure-gap}. That step is an identity in
  $W_{2,\Gamma}^2$ and holds for any reference cost placed there;
  equation~\ref{eq:random-exposure-excess-nonneg} is the separate consequence
  of that particular reference cost being the infimum of the same expectation
  over the same coupling set. The ceiling is attained rather than merely
  approached because
  the quadratic cost is lower semicontinuous and the coupling set is weakly
  compact, so an optimal plan exists and realizes $\bar\kappa_{ij}$
  \citep[Ch.~1]{panaretos_invitation_2020}.
\end{proof}

Conditional on the maintained exposure laws, the theorem restricts feasible
systematic covariance rather than selecting the realized coupling or
determining expected returns.
Holding $v_i$ and $v_j$ fixed, a larger latent exposure distance lowers the
greatest systematic covariance that the two laws can attain.
The transport excess $\Delta_\pi$ measures how much the realized joint
arrangement exceeds the minimum transport cost and, equivalently, how far its
covariance falls below the ceiling.
This is a covariance-envelope result, not a no-arbitrage pricing model.

To characterize the entire feasible covariance set rather than only its
ceiling, the same argument is applied to the reflected law $\check P_j:=(-I)_\#P_j$.

\begin{theorem}[Reflected lower endpoint]
  \label{thm:random-exposure-lower-endpoint}
  Under \Cref{ass:random-common-map}, the least systematic covariance
  attainable over the coupling set is
  \begin{equation}\label{eq:random-exposure-floor}
    \underline\kappa_{ij}
    =\frac12\left(W_{2,\Gamma}^{2}(P_i,\check P_j)-v_i-v_j\right).
  \end{equation}
\end{theorem}

\begin{proof}
  Define the reflection map
  \[
    R(z_i,z_j):=(z_i,-z_j).
  \]
  For every coupling $\pi$ of $(P_i,P_j)$, the pushforward
  $R_\#\pi$ is a coupling of $(P_i,\check P_j)$, where
  $\check P_j:=(-I)_\#P_j$. Since $R$ is involutive, this operation bijects
  the couplings of $(P_i,P_j)$ with those of $(P_i,\check P_j)$. Reflection
  leaves the two second moments unchanged and negates the covariance:
  \[
    \E_{R_\#\pi}\inner{Z_i,Z_j}_{\mathcal H}
    =-\E_\pi\inner{Z_i,Z_j}_{\mathcal H}.
  \]
  Apply \Cref{thm:random-exposure-envelope} to the reflected pair. Its
  greatest covariance is
  $\frac12(v_i+v_j-W_{2,\Gamma}^{2}(P_i,\check P_j))$; negating this value
  gives the least covariance for the original pair, which is
  equation~\ref{eq:random-exposure-floor}.
\end{proof}

Economically, the reflected endpoint describes the weakest systematic
covariance compatible with the two marginal exposure laws under their most
oppositely aligned feasible coupling.
It completes the envelope; it does not predict that realized co-movement will
be negative.

Together with \Cref{thm:random-exposure-envelope}, this brackets systematic
covariance in $[\underline\kappa_{ij},\bar\kappa_{ij}]$.
If $\pi^-$ and $\pi^+$ attain the lower and upper endpoints, respectively,
then for every $\lambda\in[0,1]$ the mixture
\[
  \pi_\lambda:=(1-\lambda)\pi^-+\lambda\pi^+
\]
preserves both marginals and, by linearity of covariance, attains
$(1-\lambda)\underline\kappa_{ij}+\lambda\bar\kappa_{ij}$.
Thus every interior covariance is attainable.
The marginal laws therefore bound covariance but do not generally determine
its sign; the restriction that survives empirically is the nonnegative
transport excess in equation~\ref{eq:random-exposure-excess-nonneg}.

The covariance envelope still depends on the latent distance
$W_{2,\Gamma}(P_i,P_j)$, which is not observed in the application.
The next step is therefore to use \Cref{ass:random-common-map} to transfer
bounds from characteristic-side distance, later proxied by article
embeddings, to exposure-side distance.
The finite-support result first shows how this bridge operates for
characteristic clouds of the form used empirically.
It pairs each characteristic atom with an indexed transmission draw so that
the characteristic and exposure clouds retain the same atom masses.
Because additive slack changes squared distances through a cross-term, the
comparison also needs a maintained support-diameter bound.

\begin{proposition}[Finite-support transfer]\label{prop:random-w2-transfer}
  Write each finite characteristic cloud as
  \[
    \widehat C_i=\sum_a p_{ia}\delta_{x_{ia}},
    \qquad
    \widehat C_j=\sum_b p_{jb}\delta_{x_{jb}}.
  \]
  Fix one transmission draw $u_{ia}$ for each characteristic atom and define
  the indexed exposure clouds by
  \[
    \widehat P_i:=\sum_a p_{ia}\delta_{T(x_{ia},u_{ia})},
    \qquad
    \widehat P_j:=\sum_b p_{jb}\delta_{T(x_{jb},u_{jb})}.
  \]
  $G_{ij}:=W_{2,d_\Omega}(\widehat C_i,\widehat C_j)$ and
  $\tau_{ij}:=\tau_i+\tau_j$. Let $D_{ij}$ be any maintained bound on the
  risk-coordinate support diameter, that is on
  $\norm{t_\Gamma(x)-t_\Gamma(y)}$ and on $\norm{Z_i-Z_j}$ across all pairs of
  indexed support points. Then, under \Cref{ass:random-common-map},
  \begin{align}
    L^{-2}G_{ij}^{2}-2\tau_{ij}D_{ij}-\tau_{ij}^{2}
    &\le W_{2,\Gamma}^{2}(\widehat P_i,\widehat P_j)
    \label{eq:random-w2-lower-transfer}\\
    W_{2,\Gamma}^{2}(\widehat P_i,\widehat P_j)
    &\le L^{2}G_{ij}^{2}+2\tau_{ij}D_{ij}+\tau_{ij}^{2}.
    \label{eq:random-w2-upper-transfer}
  \end{align}
\end{proposition}

\begin{proof}
  Let $M=(m_{ab})$ be a transport matrix. Because the indexed characteristic
  and exposure clouds use the same atom masses, every nonnegative matrix with
  row sums $p_{ia}$ and column sums $p_{jb}$ is feasible for both the
  characteristic and exposure problems. Thus both Wasserstein costs minimize
  over the same coupling-matrix polytope.
  Write $Y_{ia}:=t_\Gamma(x_{ia})$ and
  $Z_{ia}:=\Gamma^{1/2}T(x_{ia},u_{ia})$, and define $Y_{jb},Z_{jb}$
  analogously. The indexed slack bounds imply
  $\norm{Z_{ia}-Y_{ia}}\le\tau_i$ and
  $\norm{Z_{jb}-Y_{jb}}\le\tau_j$.
  For each entry $(a,b)$, the bi-Lipschitz inequality gives
  \[
    L^{-2}d_\Omega(x_{ia},x_{jb})^2
    \le \norm{t_\Gamma(x_{ia})-t_\Gamma(x_{jb})}^2
    \le L^2d_\Omega(x_{ia},x_{jb})^2.
  \]
  The indexed slack vectors have norms at most $\tau_i$ and $\tau_j$, so the
  triangle inequality gives, entry by entry,
  \[
    \left|
    \norm{Z_{ia}-Z_{jb}}^2
    -\norm{t_\Gamma(x_{ia})-t_\Gamma(x_{jb})}^2
    \right|
    \le 2\tau_{ij}D_{ij}+\tau_{ij}^2.
  \]
  Combining the two entrywise displays, average against an arbitrary feasible
  $M$. Because every admissible $M$ is available in both problems, minimizing
  the lower and upper comparisons over the shared polytope gives
  equations~\ref{eq:random-w2-lower-transfer}--\ref{eq:random-w2-upper-transfer}.
\end{proof}

The finite-support proposition shows why matching atom masses matter: the
same transport weights can be evaluated before and after transmission, so
bounded distortion and slack carry characteristic-cloud distance into an
exposure-distance interval.
Economically, the interval is conditional on the declared values of $L$,
$\tau_{ij}$, and $D_{ij}$; permitting more transmission distortion or a
larger support bound weakens what a finite characteristic cloud can imply
about its indexed exposure cloud.

For the full characteristic and exposure laws, the empirical bridge needs a
metric-level bracket that does not depend on a finite support diameter.
The following corollary supplies that population statement by comparing each
exposure law with the pushforward of its characteristic law through the common carrier.

\begin{corollary}[Metric transfer bracket]\label{cor:random-w2-metric}
  Under \Cref{ass:random-common-map}, with
  $G_{ij}:=W_{2,d_\Omega}(C_i,C_j)$, $\tau_{ij}:=\tau_i+\tau_j$, and
  $(x)_+:=\max\{x,0\}$,
  \begin{align}
    \left(L^{-1}G_{ij}-\tau_{ij}\right)_+
    &\le W_{2,\Gamma}(P_i,P_j)
    \label{eq:random-w2-lower-metric}\\
    W_{2,\Gamma}(P_i,P_j)
    &\le LG_{ij}+\tau_{ij}.
    \label{eq:random-w2-upper-metric}
  \end{align}
\end{corollary}

\begin{proof}
  Write $W_{2,\Gamma}(\mu,\nu):=W_2\bigl((\Gamma^{1/2})_\#\mu,
  (\Gamma^{1/2})_\#\nu\bigr)$, so that every leg below is measured after the
  same factor weighting and $(\Gamma^{1/2})_\#(t_\#C_i)=(t_\Gamma)_\#C_i$.
  Coupling $B_i$ with $t(X_i)$ through the shared draw $X_i$ gives
  $W_{2,\Gamma}^2(P_i,t_\#C_i)
  \le\E\norm{\Gamma^{1/2}\varepsilon_i(X_i,U_i)}^2\le\tau_i^2$, and pushing an
  optimal plan for $(C_i,C_j)$ forward through $t_\Gamma$ gives
  $W_{2,\Gamma}(t_\#C_i,t_\#C_j)\le LG_{ij}$ by the upper
  half of equation~\ref{eq:random-bilipschitz}. The $W_2$ triangle inequality
  on $\mathcal P_2(\mathcal H)$ chains the three legs into
  equation~\ref{eq:random-w2-upper-metric}.
  For the lower endpoint, let $\gamma$ be an arbitrary coupling of
  $(t_\Gamma)_\#C_i$ and $(t_\Gamma)_\#C_j$. Its marginals are supported on
  $t_\Gamma(\Omega)$, so $\gamma$ is supported on
  $t_\Gamma(\Omega)^2$. The Borel inverse on this image therefore defines
  \[
    \widetilde\gamma
    :=\bigl((t_\Gamma)^{-1},(t_\Gamma)^{-1}\bigr)_\#\gamma.
  \]
  Its marginals are $C_i$ and $C_j$. The lower half of
  equation~\ref{eq:random-bilipschitz} gives, after integration,
  \[
    G_{ij}^{2}
    \le L^{2}\int
    \norm{t_\Gamma(x)-t_\Gamma(y)}^{2}\,d\widetilde\gamma(x,y).
  \]
  Taking infima over $\gamma$ yields
  $W_{2,\Gamma}(t_\#C_i,t_\#C_j)\ge L^{-1}G_{ij}$. The same
  triangle inequality subtracts the two slack legs, and the positive part is
  taken because a Wasserstein distance is nonnegative.
\end{proof}

The bracket makes the role of the transmission assumptions transparent.
When $L=1$ and $\tau_i=\tau_j=0$, characteristic distance is preserved in
factor-risk coordinates; as multiplicative distortion or additive slack
increases, the lower endpoint falls and the upper endpoint rises.
The result therefore transfers distance conditionally rather than identifying
the latent exposure map.

To pass this exposure-distance bracket into the covariance ceiling, write
\begin{equation}\label{eq:random-w2-bracket}
  w^-_{ij}:=\left(L^{-1}W_{2,d_\Omega}(C_i,C_j)-\tau_i-\tau_j\right)_+,
  \qquad
  w^+_{ij}:=LW_{2,d_\Omega}(C_i,C_j)+\tau_i+\tau_j,
\end{equation}
and substitute \Cref{cor:random-w2-metric} into
equation~\ref{eq:random-covariance-envelope} to obtain the conditional
interval implied by characteristic-side distance and the maintained scales:
\begin{equation}\label{eq:random-kappa-bracket}
  \frac12\left(v_i+v_j-(w^+_{ij})^2\right)
  \le \bar\kappa_{ij}
  \le \frac12\left(v_i+v_j-(w^-_{ij})^2\right).
\end{equation}
The interval $[w^-_{ij},w^+_{ij}]$ is the range of latent exposure distances
allowed by the maintained transmission inputs.
Its lower endpoint has the direct economic implication used below: a larger
$w^-_{ij}$ tightens the covariance ceiling because characteristics that
cannot be brought close under any admissible transmission leave less room for
shared systematic exposure.
The next section evaluates this restriction using article distance as a proxy
for characteristic distance and treats $(L,\tau)$ as declared sensitivity
inputs rather than estimated transmission parameters.
